\section{Техническое задание}
\subsection{Основание для разработки}
Основанием для разработки является задание на выпускную квалификационную работу бакалавра «Онлайн-площадка для сдачи вещей в аренду. Разработка серверной части».
\subsection{Цель и назначение разработки}
Целью выпускной квалификационной работы является проектирование, разработка и тестирование серверной части веб-приложения, обеспечивающей хранение данных, бизнес-логику и предоставление REST API для онлайн-площадки аренды вещей. Серверная часть предоставляет арендодателям возможность публиковать объявления, а арендаторам — находить, бронировать и оплачивать аренду с гарантиями безопасности сделки.
Задачами данной разработки являются:
\begin{itemize}
\item изучение и анализ предметной области;
\item проектирование архитектуры серверной части и базы данных;
\item разработка REST API веб-приложения;
\item интеграция с  сервисом DaData;
\item тестирование разработанного API.
\end{itemize}
\subsection{Требования к серверной части веб-приложения}
\subsubsection{Функциональные требования к серверной части}
На основании описанных бизнес-процессов и анализа предметной области сформирован перечень функциональных требований к API и серверной части, представленный в таблице \ref{analiz:table}.
\begin{xltabular}{\textwidth}{|c|X|X|}
 \caption{Функциональные требования к серверной части \label{analiz:table}}\\ \hline
  № & \centrow Функциональное требование & \centrow Описание \\ \hline
 \endfirsthead
 \continuecaption{Продолжение таблицы \ref{analiz:table}}
 ~ № & \centrow Функциональное требование & \centrow Описание \\ \hline 
 \finishhead
 \leftrow 1 & Регистрация и аутентификация & \leftrow Создание учётной записи с подтверждением email через шестизначный код; выдача пары JWT-токенов при входе; восстановление пароля по коду из письма. \\ \hline
 \leftrow 2 & Верификация пользователей & \leftrow Серверная валидация реквизитов в зависимости от типа пользователя: паспорт и ИНН для физических лиц, ИНН и ОГРНИП для ИП, ИНН и КПП для юридических лиц; автозаполнение реквизитов по ИНН через DaData. \\ \hline
 \leftrow 3 & Управление профилем & \leftrow API-endpoints чтения и обновления персональных данных; пересчёт признака заполненности профиля; смена пароля. \\ \hline
 \leftrow 4 & Размещение объявления & \leftrow Создание записи объявления с атрибутами: наименование, категория, описание, цена за день, залог; назначение владельца; отправка на модерацию. \\ \hline
 \leftrow 5 & Управление объявлениями & \leftrow Endpoints редактирования и удаления собственных объявлений с проверкой прав владельца; загрузка и удаление изображений с сохранением в файловом хранилище. \\ \hline
 \leftrow 6 & Фильтрация и поиск & \leftrow Серверная фильтрация объявлений по категории и периоду доступности с исключением занятых дат из активных бронирований; полнотекстовый поиск по наименованию и описанию. \\ \hline
 \leftrow 7 & Карточка объявления & \leftrow Предоставление полных данных объявления: характеристики, изображения, рейтинг, отзывы, занятые периоды аренды, вычисленные из активных бронирований. \\ \hline
 \leftrow 8 & Онлайн-бронирование & \leftrow Создание бронирования с проверкой доступности дат; автоматический расчёт стоимости на сервере с учётом наценки платформы 20\% и сервисного сбора 5\%; управление статусами бронирования. \\ \hline
 \leftrow 9 & Онлайн-оплата & \leftrow Формирование демонстрационной СБП-сессии с реквизитами и QR-кодом; обновление статуса оплаты после подтверждения; обработка запроса на возврат средств через endpoint. \\ \hline
 \leftrow 10 & Система отзывов & \leftrow Создание отзыва по завершённому или оплаченному бронированию; автоматический пересчёт среднего рейтинга объявления и пользователя после каждого нового отзыва. \\ \hline
 \leftrow 11 & Аналитика арендодателя & \leftrow Вычисление статистики на сервере: количество бронирований, чистая прибыль без учёта комиссий и залога, история сделок за период. \\ \hline
 \leftrow 12 & Административная панель & \leftrow Модерация объявлений и пользователей, управление категориями, просмотр транзакций и заявок на модерацию через встроенную панель администратора Django. \\ \hline
\end{xltabular}
\subsubsection{Требования к программному интерфейсу}
Серверная часть должна предоставлять следующие группы API-endpoints:
\begin{itemize}
\item модуль пользователей: регистрация, подтверждение email, аутентификация, управление профилем, смена и восстановление пароля, автозаполнение реквизитов по ИНН;
\item модуль объявлений: создание, редактирование, удаление, поиск, фильтрация, загрузка и удаление изображений, управление статусами модерации;
\item модуль бронирований: создание заявки, подтверждение и отмена, проведение оплаты, подтверждение передачи и возврата товара с фотофиксацией, запрос и обработка возврата средств.
\end{itemize}
\subsection{Моделирование вариантов использования}
Для серверной части разработана модель вариантов использования на языке UML, отражающая взаимодействие пользователей с программным интерфейсом системы. На основании анализа предметной области реализованы следующие прецеденты:
\begin{enumerate}
\item Регистрация с подтверждением email.
\item Аутентификация с JWT-токен.
\item Восстановление пароля.
\item Просмотр профиля.
\item Обновление профиля.
\item Автозаполнение реквизитов по ИНН.
\item Просмотр каталога объявлений.
\item Фильтрация и поиск объявлений.
\item Просмотр карточки объявления.
\item Создание объявления.
\item Редактирование объявления.
\item Удаление объявления.
\item Загрузка изображений.
\item Удаление изображений.
\item Модерация объявлений.
\item Создание категории.
\item Удаление категории.
\item Блокировка пользователя.
\item Разблокировка пользователя.
\item Блокировка объявления.
\item Разблокировка объявления.
\item Создание бронирования.
\item Просмотр бронирований.
\item Подтверждение бронирования.
\item Отмена бронирования.
\item Завершение аренды.
\item Оплата через СБП.
\item Подтверждение получения товара.
\item Подтверждение возврата товара.
\item Возврат залога.
\item Возврат средств через СБП.
\item Создание отзыва.
\end{enumerate}

Разрабатываемое веб-приложение имеет 3 группы пользователей с разными правами: гость, пользователь, администратор.

Гостю должны быть доступны следующие функции:
\begin{enumerate}
\item Регистрация с подтверждением email.
\item Аутентификация с JWT-токен.
\item Восстановление пароля.
\item Автозаполнение реквизитов по ИНН.
\item Просмотр каталога объявлений.
\item Фильтрация и поиск объявлений.
\item Просмотр карточки объявления.
\end{enumerate}
На рисунке \ref{diaprecguets:image} изображена диаграмма прецедентов для гостя.
\begin{figure}[H]
\center{\includegraphics[width=1\linewidth]{diaprecguets}}
\caption{Диаграмма прецедентов для гостя}
\label{diaprecguets:image}
\end{figure}

Пользователю должны быть доступны следующие функции:
\begin{enumerate}
\item Просмотр профиля.
\item Обновление профиля.
\item Автозаполнение реквизитов по ИНН.
\item Просмотр каталога объявлений.
\item Фильтрация и поиск объявлений.
\item Просмотр карточки объявления.
\item Создание объявления.
\item Редактирование объявления.
\item Удаление объявления.
\item Загрузка изображений.
\item Удаление изображений.
\item Создание бронирования.
\item Просмотр бронирований.
\item Подтверждение бронирования.
\item Отмена бронирования.
\item Завершение аренды.
\item Оплата через СБП.
\item Подтверждение получения товара.
\item Подтверждение возврата товара.
\item Возврат залога.
\item Возврат средств через СБП.
\item Создание отзыва.
\end{enumerate}
На рисунке \ref{diaprecuser:image} изображена диаграмма прецедентов для пользователя.
\begin{figure}[H]
\center{\includegraphics[width=1\linewidth]{diaprecuser}}
\caption{Диаграмма прецедентов для пользователя}
\label{diaprecuser:image}
\end{figure}

Администратору должны быть доступны следующие функции:
\begin{enumerate}
\item Просмотр профиля.
\item Просмотр каталога объявлений.
\item Фильтрация и поиск объявлений.
\item Просмотр карточки объявления.
\item Просмотр бронирований.
\item Модерация объявлений.
\item Создание категории.
\item Удаление категории.
\item Блокировка пользователя.
\item Разблокировка пользователя.
\item Блокировка объявления.
\item Разблокировка объявления.
\end{enumerate}
На рисунке \ref{diaprecadmin:image} изображена диаграмма прецедентов для администратора.
\begin{figure}[H]
\center{\includegraphics[width=1\linewidth]{diaprecadmin}}
\caption{Диаграмма прецедентов для администратора}
\label{diaprecadmin:image}
\end{figure}

\subsubsection{Регистрация с подтверждением email}
Заинтересованные лица и их требования: пользователь желает создать учётную запись для получения доступа к функциям сервиса аренды вещей.
Предусловие: пользователь обращается к API с корректными данными; указанный email не зарегистрирован в системе.
Постусловие: в базе данных создана новая учётная запись; email подтверждён.
Основной успешный сценарий:
\begin{enumerate}
\item Клиент отправляет POST-запрос с адресом электронной почты.
\item Сервер проверяет уникальность email в базе данных.
\item Сервер генерирует шестизначный код подтверждения и сохраняет его в кэше с таймаутом 15 минут.
\item Сервер отправляет код на указанный email через SMTP.
\item Клиент отправляет POST-запрос с именем пользователя, email, паролем, типом аккаунта, реквизитами и кодом подтверждения из письма.
\item Сервер проверяет соответствие кода сохранённому значению и его срок действия.
\item Сервер валидирует реквизиты в соответствии с типом пользователя: для физического лица — ФИО и паспортные данные; для ИП — наименование, ИНН и ОГРНИП; для юридического лица — наименование, ИНН и КПП.
\item Сервер создаёт учётную запись, сохраняет её в базе данных и удаляет использованный код из кэша.
\item Сервер возвращает клиенту данные созданного пользователя.
\end{enumerate}

\subsubsection{Аутентификация и JWT-токен}
Заинтересованные лица и их требования: пользователь желает войти в систему для получения доступа к защищённым ресурсам API.
Предусловие: учётная запись пользователя существует в базе данных.
Постусловие: пользователь получает пару JWT-токенов для последующей авторизации запросов.
Основной успешный сценарий:
\begin{enumerate}
\item Пользователь направляет серверу логин и пароль.
\item Сервер проверяет соответствие переданных учётных данных записи в базе данных.
\item Сервер формирует и возвращает access-токен с ограниченным сроком действия и refresh-токен для его обновления.
\item При последующих обращениях к защищённым endpoints клиент передаёт access-токен в заголовке запроса.
\item При истечении срока действия access-токена клиент использует refresh-токен для получения нового.
\end{enumerate}

\subsubsection{Восстановление пароля}
Заинтересованные лица и их требования: пользователь забыл пароль и желает восстановить доступ к учётной записи.
Предусловие: учётная запись с указанным email существует в базе данных.
Постусловие: пароль пользователя изменён на новый; использованный код удалён из кэша.
Основной успешный сценарий:
\begin{enumerate}
\item Пользователь направляет серверу запрос с адресом электронной почты.
\item Сервер проверяет наличие учётной записи с указанным email.
\item Сервер генерирует шестизначный код и сохраняет его в кэше с таймаутом 15 минут.
\item Сервер отправляет код на email пользователя.
\item Пользователь направляет серверу запрос с email, кодом из письма и новым паролем.
\item Сервер проверяет соответствие кода сохранённому значению и его срок действия.
\item Сервер устанавливает новый пароль и удаляет использованный код из кэша.
\end{enumerate}

\subsubsection{Просмотр профиля}
Заинтересованные лица и их требования: авторизованный пользователь желает получить актуальные данные своей учётной записи.
Предусловие: пользователь авторизован в системе.
Постусловие: сервер возвращает данные профиля текущего пользователя.
Основной успешный сценарий:
\begin{enumerate}
\item Пользователь запрашивает данные своего профиля.
\item Сервер извлекает идентификатор пользователя из токена и находит его запись в базе данных.
\item Сервер сериализует данные профиля: имя пользователя, адрес электронной почты, тип аккаунта, реквизиты и признак заполненности профиля.
\item Сервер возвращает сериализованные данные в формате JSON.
\end{enumerate}

\subsubsection{Обновление профиля}
Заинтересованные лица и их требования: пользователь желает обновить данные своей учётной записи.
Предусловие: пользователь авторизован.
Постусловие: данные профиля обновлены в базе данных; признак заполненности профиля пересчитан.
Основной успешный сценарий:
\begin{enumerate}
\item Пользователь направляет серверу обновляемые поля профиля.
\item Сервер валидирует переданные поля в соответствии с типом пользователя.
\item Сервер сохраняет изменения и пересчитывает признак заполненности профиля.
\item Сервер возвращает обновлённые данные профиля.
\end{enumerate}

\subsubsection{Автозаполнение реквизитов по ИНН}
Заинтересованные лица и их требования: ИП и юридические лица хотят автоматически заполнить реквизиты по ИНН без ручного ввода.
Предусловие: пользователь знает свой ИНН; внешний сервис DaData доступен.
Постусловие: сервер возвращает реквизиты, соответствующие переданному ИНН.
Основной успешный сценарий:
\begin{enumerate}
\item Пользователь направляет серверу запрос с ИНН.
\item Сервер обращается к внешнему сервису DaData и определяет тип субъекта по длине ИНН: 10 цифр — юридическое лицо, 12 цифр — ИП.
\item Сервер возвращает наименование организации и соответствующие реквизиты: для юридического лица — ИНН и КПП, для ИП — ИНН и ОГРНИП.
\item Пользователь использует полученные данные для заполнения профиля.
\end{enumerate}

\subsubsection{Просмотр каталога объявлений}
Заинтересованные лица и их требования: пользователь желает получить список вещей, доступных для аренды.
Предусловие: запрос может выполнять как авторизованный, так и неавторизованный пользователь.
Постусловие: сервер возвращает список опубликованных объявлений.
Основной успешный сценарий:
\begin{enumerate}
\item Пользователь запрашивает список объявлений.
\item Сервер формирует запрос к базе данных: для неавторизованного пользователя выбирает только объявления со статусом «Опубликовано», для авторизованного — также включает его собственные объявления в любом статусе.
\item Сервер сериализует данные: наименование, описание, цену за день, залог, категорию, средний рейтинг, количество отзывов и ссылки на изображения.
\item Сервер возвращает постраничный список объявлений в формате JSON.
\end{enumerate}

\subsubsection{Фильтрация и поиск объявлений}
Заинтересованные лица и их требования: пользователь желает получить только те объявления, которые соответствуют заданным критериям.
Предусловие: пользователь передаёт параметры фильтрации или поиска в запросе.
Постусловие: сервер возвращает отфильтрованный список объявлений.
Основной успешный сценарий:
\begin{enumerate}
\item Пользователь передаёт параметры в запросе: категорию, текст поиска, порядок сортировки или минимальный рейтинг.
\item Сервер применяет фильтр по идентификатору категории и полнотекстовый поиск по названию и описанию.
\item Сервер применяет указанную сортировку по цене, дате создания или рейтингу.
\item Сервер возвращает отфильтрованный постраничный список объявлений.
\end{enumerate}

\subsubsection{Просмотр карточки объявления}
Заинтересованные лица и их требования: пользователь желает получить полные данные о вещи перед оформлением аренды.
Предусловие: объявление существует в базе данных.
Постусловие: сервер возвращает полный набор данных по объявлению, включая рейтинг, отзывы и занятые периоды.
Основной успешный сценарий:
\begin{enumerate}
\item Пользователь запрашивает данные конкретного объявления по его идентификатору.
\item Сервер выбирает запись объявления вместе с данными владельца, категории и изображений.
\item Сервер прикрепляет список отзывов и средний рейтинг объявления и его владельца.
\item Сервер формирует список занятых периодов на основе активных бронирований со статусами «Рассматривается» и «Подтверждено».
\item Сервер возвращает полный набор данных в формате JSON.
\end{enumerate}

\subsubsection{Создание объявления}
Заинтересованные лица и их требования: авторизованный пользователь с заполненным профилем желает разместить вещь для аренды.
Предусловие: пользователь авторизован; признак заполненности профиля имеет значение «заполнен».
Постусловие: объявление создано в базе данных со статусом «На модерации»; заявка на модерацию зафиксирована.
Основной успешный сценарий:
\begin{enumerate}
\item Пользователь направляет серверу POST-запрос с данными объявления: наименование, описание, цена за день, залог, категория.
\item Сервер проверяет заполненность профиля пользователя; при незаполненном профиле возвращает ошибку.
\item Сервер создаёт запись объявления со статусом «На модерации» и назначает пользователя владельцем.
\item Сервер создаёт заявку на модерацию с типом действия «Создание».
\item Сервер возвращает созданное объявление с присвоенным идентификатором.
\end{enumerate}

\subsubsection{Редактирование объявления}
Заинтересованные лица и их требования: владелец желает изменить данные ранее созданного объявления.
Предусловие: пользователь авторизован, является владельцем объявления и имеет заполненный профиль.
Постусловие: данные объявления обновлены в базе данных; создана заявка на повторную модерацию.
Основной успешный сценарий:
\begin{enumerate}
\item Пользователь направляет серверу запрос с изменёнными полями объявления.
\item Сервер проверяет, что текущий пользователь является владельцем данного объявления; в противном случае возвращает ошибку доступа.
\item Сервер сохраняет обновлённые данные, переводит статус объявления в «На модерации» и создаёт заявку на модерацию с типом действия «Редактирование».
\item Сервер возвращает обновлённую карточку объявления.
\end{enumerate}

\subsubsection{Удаление объявления}
Заинтересованные лица и их требования: владелец желает удалить объявление из системы.
Предусловие: пользователь авторизован и является владельцем объявления.
Постусловие: объявление и связанные изображения удалены из базы данных и файлового хранилища.
Основной успешный сценарий:
\begin{enumerate}
\item Пользователь направляет серверу DELETE-запрос на удаление объявления.
\item Сервер проверяет, что текущий пользователь является владельцем данного объявления; в противном случае возвращает ошибку доступа.
\item Сервер удаляет объявление и связанные с ним изображения из базы данных и файлового хранилища.
\item Сервер возвращает подтверждение успешного удаления.
\end{enumerate}

\subsubsection{Загрузка изображений}
Заинтересованные лица и их требования: владелец желает прикрепить фотографии вещи к объявлению.
Предусловие: пользователь авторизован и является владельцем объявления.
Постусловие: изображение сохранено в файловом хранилище и связано с объявлением в базе данных.
Основной успешный сценарий:
\begin{enumerate}
\item Пользователь направляет серверу POST-запрос с файлом изображения и идентификатором объявления в формате multipart/form-data.
\item Сервер проверяет, что текущий пользователь является владельцем указанного объявления; в противном случае возвращает ошибку доступа.
\item Сервер сохраняет файл в файловом хранилище.
\item Сервер создаёт запись, связывающую изображение с объявлением.
\item Сервер возвращает идентификатор и ссылку на загруженное изображение.
\end{enumerate}

\subsubsection{Удаление изображений}
Заинтересованные лица и их требования: владелец желает удалить изображение из своего объявления.
Предусловие: пользователь авторизован и является владельцем объявления; изображение привязано к данному объявлению.
Постусловие: изображение удалено из базы данных и файлового хранилища.
Основной успешный сценарий:
\begin{enumerate}
\item Пользователь направляет серверу DELETE-запрос на удаление конкретного изображения.
\item Сервер проверяет, что изображение принадлежит объявлению текущего пользователя; в противном случае возвращает ошибку доступа.
\item Сервер удаляет файл из хранилища и запись из базы данных.
\item Сервер возвращает подтверждение успешного удаления.
\end{enumerate}

\subsubsection{Модерация объявлений}
Заинтересованные лица и их требования: администратор желает проверить поступившие заявки на публикацию объявлений и принять по ним решение.
Предусловие: в системе существуют заявки на модерацию со статусом «Ожидает рассмотрения»; пользователь обладает правами суперпользователя.
Постусловие: статус заявки и связанного объявления обновлены; владелец уведомлён о результате.
Основной успешный сценарий:
\begin{enumerate}
\item Администратор запрашивает список заявок на модерацию, при необходимости фильтруя по статусу.
\item Сервер возвращает заявки с полными данными объявления, включая изображения, описание и реквизиты владельца.
\item Администратор принимает решение об одобрении: сервер переводит статус объявления в «Опубликовано» и статус заявки в «Одобрено».
\item При отклонении администратор направляет запрос с указанием причины: сервер переводит статус объявления в «Отклонено», статус заявки в «Отклонено» и сохраняет причину отказа.
\end{enumerate}

\subsubsection{Создание категории}
Заинтересованные лица и их требования: администратор желает добавить новую категорию товаров для структурирования каталога.
Предусловие: пользователь авторизован и обладает правами суперпользователя.
Постусловие: категория создана в базе данных и доступна для выбора при размещении объявлений.
Основной успешный сценарий:
\begin{enumerate}
\item Администратор направляет серверу POST-запрос с наименованием новой категории.
\item Сервер проверяет права доступа пользователя.
\item Сервер создаёт запись категории в базе данных.
\item Сервер возвращает созданную категорию с присвоенным идентификатором.
\end{enumerate}

\subsubsection{Удаление категории}
Заинтересованные лица и их требования: администратор желает удалить устаревшую или ошибочно созданную категорию.
Предусловие: пользователь авторизован и обладает правами суперпользователя; категория существует в базе данных.
Постусловие: категория удалена из базы данных; связанные объявления остаются без категории.
Основной успешный сценарий:
\begin{enumerate}
\item Администратор направляет серверу DELETE-запрос на удаление категории по её идентификатору.
\item Сервер проверяет права доступа пользователя.
\item Сервер удаляет запись категории из базы данных.
\item Сервер возвращает подтверждение успешного удаления.
\end{enumerate}

\subsubsection{Блокировка пользователя}
Заинтересованные лица и их требования: администратор желает ограничить доступ недобросовестного пользователя к платформе.
Предусловие: пользователь обладает правами суперпользователя; целевая учётная запись существует в базе данных и не является суперпользователем.
Постусловие: признак активности учётной записи установлен в значение «неактивен»; пользователь лишается возможности аутентификации.
Основной успешный сценарий:
\begin{enumerate}
\item Администратор направляет серверу запрос на блокировку пользователя с указанием его идентификатора.
\item Сервер проверяет права доступа и убеждается, что целевой пользователь не является суперпользователем.
\item Сервер устанавливает признак активности учётной записи в значение «неактивен».
\item Сервер возвращает подтверждение выполненной блокировки.
\end{enumerate}

\subsubsection{Разблокировка пользователя}
Заинтересованные лица и их требования: администратор желает восстановить доступ ранее заблокированного пользователя к платформе.
Предусловие: пользователь обладает правами суперпользователя; целевая учётная запись существует в базе данных и имеет статус «неактивен».
Постусловие: признак активности учётной записи установлен в значение «активен»; пользователь восстанавливает возможность аутентификации.
Основной успешный сценарий:
\begin{enumerate}
\item Администратор направляет серверу запрос на разблокировку пользователя с указанием его идентификатора.
\item Сервер проверяет права доступа.
\item Сервер устанавливает признак активности учётной записи в значение «активен».
\item Сервер возвращает подтверждение восстановления доступа.
\end{enumerate}

\subsubsection{Блокировка объявления}
Заинтересованные лица и их требования: администратор желает скрыть объявление, нарушающее правила платформы, из публичного каталога.
Предусловие: пользователь обладает правами суперпользователя; объявление существует в базе данных.
Постусловие: статус объявления переведён в «Заблокировано»; объявление не отображается в каталоге для других пользователей.
Основной успешный сценарий:
\begin{enumerate}
\item Администратор направляет серверу запрос на блокировку объявления с указанием его идентификатора.
\item Сервер проверяет права доступа пользователя.
\item Сервер переводит статус объявления в «Заблокировано».
\item Сервер возвращает подтверждение выполненной блокировки.
\end{enumerate}

\subsubsection{Разблокировка объявления}
Заинтересованные лица и их требования: администратор желает восстановить видимость ранее заблокированного объявления в публичном каталоге.
Предусловие: пользователь обладает правами суперпользователя; объявление существует в базе данных и имеет статус «Заблокировано».
Постусловие: статус объявления переведён в «Опубликовано»; объявление снова отображается в каталоге.
Основной успешный сценарий:
\begin{enumerate}
\item Администратор направляет серверу запрос на разблокировку объявления с указанием его идентификатора.
\item Сервер проверяет права доступа пользователя.
\item Сервер переводит статус объявления в «Опубликовано».
\item Сервер возвращает подтверждение восстановления видимости объявления.
\end{enumerate}

\subsubsection{Создание бронирования}
Заинтересованные лица и их требования: авторизованный пользователь желает забронировать чужую вещь; итоговая стоимость должна рассчитываться автоматически на стороне сервера.
Предусловие: пользователь авторизован, его профиль заполнен; пользователь не является владельцем выбранного объявления.
Постусловие: бронирование создано; итоговая стоимость с разбивкой по компонентам сохранена в базе данных.
Основной успешный сценарий:
\begin{enumerate}
\item Пользователь направляет серверу POST-запрос с идентификатором объявления и датами начала и окончания аренды.
\item Сервер проверяет, что пользователь не является владельцем вещи и что его профиль заполнен.
\item Сервер проверяет отсутствие пересечения запрошенного периода с существующими активными бронированиями.
\item Сервер рассчитывает количество дней аренды и базовую стоимость.
\item Сервер применяет наценку платформы 20\% и сервисный сбор 5\%, фиксирует сумму залога из объявления.
\item Сервер сохраняет бронирование со статусом «Рассматривается» и отдельными полями: сумма аренды, сервисный сбор, сумма залога и итоговая сумма.
\item Сервер возвращает созданное бронирование с полной разбивкой стоимости.
\end{enumerate}

\subsubsection{Просмотр бронирований}
Заинтересованные лица и их требования: пользователь желает просматривать как заявки на аренду своих вещей, так и заявки, оформленные им самим.
Предусловие: пользователь авторизован.
Постусловие: сервер возвращает список бронирований с возможностью фильтрации по роли.
Основной успешный сценарий:
\begin{enumerate}
\item Пользователь запрашивает список бронирований, при необходимости указывая роль: входящие или исходящие.
\item Сервер выбирает из базы данных бронирования, в которых текущий пользователь является арендатором или владельцем вещи.
\item При наличии параметра роли сервер дополнительно фильтрует список: входящие — только заявки на собственные вещи, исходящие — только заявки пользователя.
\item Сервер сериализует данные: вещь, даты, стоимость с разбивкой по компонентам, статус бронирования, статус оплаты, статус залога.
\item Сервер возвращает постраничный список бронирований в формате JSON.
\end{enumerate}

\subsubsection{Подтверждение бронирования}
Заинтересованные лица и их требования: арендодатель желает принять входящую заявку на аренду своей вещи.
Предусловие: бронирование имеет статус «Рассматривается»; текущий пользователь является владельцем вещи.
Постусловие: статус бронирования изменён на «Подтверждено»; арендатор уведомлён.
Основной успешный сценарий:
\begin{enumerate}
\item Арендодатель направляет серверу POST-запрос на подтверждение бронирования.
\item Сервер проверяет, что текущий пользователь является владельцем данной вещи; в противном случае возвращает ошибку доступа.
\item Сервер переводит статус бронирования в «Подтверждено».
\item Сервер возвращает обновлённые данные бронирования.
\end{enumerate}

\subsubsection{Отмена бронирования}
Заинтересованные лица и их требования: арендатор или арендодатель желают отменить бронирование; арендатор может указать причину отказа.
Предусловие: бронирование существует в базе данных и находится в допустимом для отмены статусе.
Постусловие: статус бронирования изменён на «Отменено»; причина отказа сохранена при наличии.
Основной успешный сценарий:
\begin{enumerate}
\item Участник сделки направляет серверу POST-запрос на отмену бронирования.
\item Сервер проверяет, что текущий пользователь является арендатором или владельцем данного бронирования; в противном случае возвращает ошибку доступа.
\item Сервер переводит статус бронирования в «Отменено».
\item При отказе арендатора с указанием причины сервер дополнительно сохраняет текст причины отказа.
\item Сервер возвращает обновлённые данные бронирования.
\end{enumerate}

\subsubsection{Завершение аренды}
Заинтересованные лица и их требования: владелец желает завершить аренду после возврата вещи арендатором.
Предусловие: бронирование имеет статус «Подтверждено» и статус оплаты «Оплачено».
Постусловие: статус бронирования изменён на «Завершено».
Основной успешный сценарий:
\begin{enumerate}
\item Арендодатель направляет серверу POST-запрос на завершение аренды.
\item Сервер проверяет, что текущий пользователь является владельцем данной вещи.
\item Сервер проверяет, что бронирование оплачено.
\item Сервер переводит статус бронирования в «Завершено».
\item Сервер возвращает обновлённые данные бронирования.
\end{enumerate}

\subsubsection{Оплата через СБП}
Заинтересованные лица и их требования: арендатор желает оплатить подтверждённое бронирование через демонстрационный механизм СБП.
Предусловие: бронирование имеет статус «Подтверждено»; оплата ещё не произведена.
Постусловие: статус оплаты бронирования изменён на «Оплачено»; при наличии залога его статус также изменён на «Оплачен».
Основной успешный сценарий:
\begin{enumerate}
\item Арендатор направляет POST-запрос на запуск оплаты.
\item Сервер проверяет статус бронирования, его принадлежность данному арендатору и отсутствие действующей платёжной сессии.
\item Сервер формирует демонстрационную СБП-сессию: генерирует идентификатор платежа, устанавливает срок действия 15 минут и формирует реквизиты платежа (сумма, получатель, банк, QR-данные).
\item Сервер возвращает параметры платежа арендатору.
\item Арендатор подтверждает факт оплаты, направляя POST-запрос.
\item Сервер проверяет актуальность платёжной сессии; при истечении срока переводит статус в «Ошибка».
\item Сервер устанавливает статус оплаты «Оплачено», фиксирует дату оплаты и, при наличии залога, переводит его статус в «Оплачен».
\end{enumerate}

\subsubsection{Подтверждение получения товара}
Заинтересованные лица и их требования: арендатор желает зафиксировать факт получения вещи в начале аренды.
Предусловие: бронирование имеет статус «Подтверждено» и статус оплаты «Оплачено».
Постусловие: факт получения товара зафиксирован в базе данных с фотографией и временной меткой.
Основной успешный сценарий:
\begin{enumerate}
\item Арендатор направляет POST-запрос с прикреплённой фотографией в формате multipart/form-data.
\item Сервер проверяет, что текущий пользователь является арендатором данного бронирования и что оплата произведена.
\item Сервер сохраняет фотографию в файловом хранилище.
\item Сервер устанавливает признак получения товара арендатором и фиксирует дату и время подтверждения.
\item Сервер возвращает обновлённые данные бронирования.
\end{enumerate}

\subsubsection{Подтверждение возврата товара}
Заинтересованные лица и их требования: арендодатель желает зафиксировать факт возврата вещи по окончании аренды.
Предусловие: арендатор ранее подтвердил получение товара; бронирование имеет статус «Подтверждено».
Постусловие: факт возврата товара зафиксирован в базе данных с фотографией и временной меткой; бронирование автоматически переведено в статус «Завершено».
Основной успешный сценарий:
\begin{enumerate}
\item Арендодатель направляет POST-запрос с фотографией возвращённого товара в формате multipart/form-data.
\item Сервер проверяет, что текущий пользователь является владельцем данной вещи.
\item Сервер проверяет, что арендатор ранее подтвердил получение товара.
\item Сервер сохраняет фотографию в файловом хранилище.
\item Сервер устанавливает признак возврата товара и автоматически переводит бронирование в статус «Завершено».
\item Сервер возвращает обновлённые данные бронирования.
\end{enumerate}

\subsubsection{Возврат залога}
Заинтересованные лица и их требования: владелец желает вернуть залог арендатору после завершения аренды.
Предусловие: бронирование завершено; статус залога имеет значение «Оплачен».
Постусловие: статус залога изменён на «Возвращён».
Основной успешный сценарий:
\begin{enumerate}
\item Арендодатель направляет POST-запрос на возврат залога.
\item Сервер проверяет, что текущий пользователь является владельцем данной вещи.
\item Сервер проверяет, что залог был оплачен.
\item Сервер переводит статус залога в «Возвращён» и возвращает обновлённые данные бронирования.
\end{enumerate}

\subsubsection{Возврат средств через СБП}
Заинтересованные лица и их требования: арендатор желает вернуть оплаченные средства в случае невозможности исполнения договора аренды; арендодатель принимает решение об одобрении или отклонении запроса.
Предусловие: бронирование имеет статус оплаты «Оплачено»; запрос на возврат ещё не подавался.
Постусловие: при одобрении статус возврата изменён на «Выполнен»; при наличии залога его статус изменён на «Возвращён».
Основной успешный сценарий:
\begin{enumerate}
\item Арендатор направляет POST-запрос на возврат средств с указанием причины.
\item Сервер проверяет, что текущий пользователь является арендатором и что запрос на возврат ранее не подавался.
\item Сервер фиксирует запрос: сохраняет причину, дату обращения и переводит статус возврата в значение «Запрошен».
\item Арендодатель получает уведомление и направляет POST-запрос с признаком одобрения.
\item Сервер проверяет, что текущий пользователь является владельцем вещи.
\item При одобрении сервер переводит статус возврата в «Выполнен» и, при наличии оплаченного залога, его статус в «Возвращён».
\item При отклонении сервер переводит статус возврата в «Отклонён»; средства остаются у арендодателя.
\end{enumerate}

\subsubsection{Создание отзыва}
Заинтересованные лица и их требования: арендатор желает оставить оценку после завершения аренды.
Предусловие: бронирование имеет статус «Завершено» либо статус «Подтверждено» при статусе оплаты «Оплачено»; отзыв по данному бронированию ещё не создавался.
Постусловие: отзыв сохранён; средний рейтинг объявления и его владельца пересчитан.
Основной успешный сценарий:
\begin{enumerate}
\item Арендатор направляет серверу POST-запрос с оценкой от 1 до 5 и текстом комментария.
\item Сервер проверяет, что текущий пользователь является арендатором данного бронирования.
\item Сервер проверяет статус бронирования и отсутствие существующего отзыва по нему.
\item Сервер сохраняет отзыв в базе данных.
\item Сервер пересчитывает средний рейтинг объявления и его владельца на основе всех видимых отзывов.
\end{enumerate}

\subsection{Требования к оформлению документации}
Разработка программной документации и программного изделия должна производиться согласно ГОСТ 19.102-77 и ГОСТ 34.601-90. Единая система программной документации.